\chapter{Энергия вакуума, инфляция и первичные флуктуации}

\section{Энергия вакуума и пространство де Ситтера}

В плоском пространстве уравнение Фридмана имеет вид
\begin{equation}
\left(\frac{\dot A}{A}\right)^2=\frac{8\pi G}{3}\rho.
\end{equation}
Для вещества \(\rho\) убывает как \(A^{-3}\). Для излучения \(\rho\) убывает как \(A^{-4}\). Для энергии вакуума она не разрежается:
\begin{equation}
\rho=\rho_0=\text{constant}.
\end{equation}
Если Вселенная содержит только эту энергию вакуума, правая часть является постоянной. Обозначая эту постоянную через \(H^2\),
\begin{equation}
H^2=\frac{8\pi G}{3}\rho_0,
\qquad
H=\sqrt{\frac{8\pi G}{3}\rho_0}.
\end{equation}
Тогда величина Хаббла действительно постоянна:
\begin{equation}
\frac{\dot A}{A}=H.
\end{equation}
Следовательно,
\begin{equation}
\dot A=HA,
\end{equation}
а решение является экспоненциальным:
\begin{equation}
A(t)=C e^{Ht}.
\end{equation}
Расстояние между далёкими несвязанными галактиками растёт экспоненциально. Это пространство де Ситтера.

В этих координатах геометрия имеет вид
\begin{equation}
d\tau^2=dt^2-e^{2Ht}\left(dx^2+dy^2+dz^2\right).
\end{equation}
Энергия вакуума здесь положительна. В плоском пространстве иначе и быть не может, поскольку левая часть уравнения Фридмана неотрицательна. Отрицательная энергия вакуума может возникнуть только при учёте членов кривизны, а этот случай здесь не рассматривается.

Эйнштейн ввёл космологическую постоянную с другой целью: уравновесить обычное вещество и удержать Вселенную статической.

\section{Современная ускоряющаяся Вселенная}

Пространство де Ситтера — это также поздневременная асимптотика современной Вселенной. Раньше масштабный фактор следовал закону, характерному для эпохи доминирования вещества,
\begin{equation}
A(t)\propto t^{2/3}.
\end{equation}
Несколько миллиардов лет назад кривая загнулась вверх и начала расти быстрее. Это ещё не точная экспонента, потому что другие формы энергии всё ещё присутствуют, но тенденция направлена к ней.

Под ускорением здесь понимается
\begin{equation}
\ddot A>0.
\end{equation}
Это наблюдаемая особенность расширения на поздних временах.

Теперь рассмотрим горизонт. Закон Хаббла гласит
\begin{equation}
v=HD.
\end{equation}
На расстоянии, где скорость удаления становится равной скорости света,
\begin{equation}
D=\frac{c}{H}.
\end{equation}
Эквивалентно это можно переписать как \(HD=c\).\editorialnote{Это редакционное алгебраическое переформулирование произнесённого соотношения \(D=c/H\). Верхняя строка на правой доске на рисунке~\ref{fig:lecture07-horizon} читается лишь частично.}

\begin{figure}[ht]
\centering
\includegraphics[width=0.62\textwidth]{lecture_07_figure_02.png}
\caption{Соотношение космического горизонта \(D=c/H\) при расширении де Ситтера.}
\label{fig:lecture07-horizon}
\lectureframe{00:30:37}
\end{figure}

Это расстояние имеет порядок от десяти до пятнадцати миллиардов световых лет и численно близко к возрасту Вселенной, если последний перевести в расстояние с помощью \(c\).

\begin{classroomqa}
\lecturetimestamp{00:14:10}
\audiencequestion{Разве почти равенство между масштабом горизонта и возрастом Вселенной не получается просто тогда, когда касательная случайно проходит через начало координат?}
\lecturerresponse{Это было бы в точности верно для \(A(t)\propto t\). В общем случае это неверно. В пространстве де Ситтера \(H\) остаётся фиксированным, тогда как возраст продолжает расти, так что нынешнее почти равенство не является математическим тождеством.}
\end{classroomqa}

\section{Что означает энергия вакуума}

Энергия вакуума описывается здесь на языке обычной квантовой механики. Для гармонического осциллятора энергия схематически равна
\begin{equation}
E\sim p^2+x^2.
\end{equation}
В классике основное состояние — это точка равновесия,
\begin{equation}
x=0,\qquad p=0,\qquad E=0.
\end{equation}
Квантовая механика не допускает, чтобы и \(x\), и \(p\) были точно зафиксированы равными нулю. Принцип неопределённости вынуждает флуктуации, так что энергия основного состояния положительна. Для осциллятора получается
\begin{equation}
E_0=\frac{1}{2}\hbar\omega.
\end{equation}

То же соображение затем переносится на поля. Вселенная полна гармонических осцилляторов: это колебания электромагнитного поля и точно так же любого другого поля. У каждого поля есть энергия нулевых колебаний. Для электромагнитного поля плотность энергии схематически равна
\begin{equation}
u_{\mathrm{EM}}\sim E^2+B^2.
\end{equation}
Даже в отсутствие реальных фотонов остаются неустранимые нулевые колебания. Это одно из объяснений энергии вакуума. Единственное свойство, которое здесь нужно, состоит в том, что эта энергия присутствует и не разрежается при расширении Вселенной.

\begin{classroomqa}
\lecturetimestamp{00:16:15}
\audiencequestion{Что представляет собой эта энергия вакуума и почему вообще что-то должно остаться, если удалить вещество и излучение?}
\lecturerresponse{Потому что вакуум — это основное состояние квантовых полей, а квантовое основное состояние не обязано иметь нулевую энергию. Гармонические осцилляторы имеют положительную энергию основного состояния, а поля ведут себя как совокупности осцилляторов. Поэтому даже в отсутствие реальных частиц остаётся энергия нулевых колебаний.}
\end{classroomqa}

Это не трёхградусное фоновое излучение. Микроволновый фон — обычное излучение. Его плотность энергии уменьшается, потому что фотоны разрежаются и каждый фотон испытывает красное смещение. Энергия вакуума — это остаточная энергия нулевых колебаний, которая остаётся после удаления реальных фотонов.

\section{Что означает \(H^{-1}\)}

Теперь вернёмся к связи между обратной величиной Хаббла и возрастом Вселенной. Если
\begin{equation}
A(t)\propto t,
\end{equation}
то
\begin{equation}
H=\frac{\dot A}{A}=\frac{1}{t}.
\end{equation}
Более общо, если
\begin{equation}
A(t)\propto t^p,
\end{equation}
то
\begin{equation}
H=\frac{p}{t}.
\end{equation}
Для случая доминирования вещества \(p=2/3\),
\begin{equation}
H=\frac{2}{3t}.
\end{equation}
Следовательно, в космологиях со степенным законом \(H^{-1}\) имеет тот же порядок величины, что и возраст, но не совпадает с ним в точности.

Пространство де Ситтера отличается. Там
\begin{equation}
A(t)=C e^{\gamma t},
\end{equation}
и, следовательно,
\begin{equation}
H=\frac{\dot A}{A}=\gamma,
\end{equation}
— подлинная постоянная. Здесь \(H^{-1}\) — не возраст Вселенной. Это постоянная времени экспоненциального роста: время, за которое масштабный фактор умножается на \(e\). В устном обсуждении несколько раз говорится ``удваивается'', а затем это исправляется; точная формулировка — умножение на \(e\). Сегодня эта постоянная времени имеет порядок десяти миллиардов лет.

\begin{classroomqa}
\lecturetimestamp{00:23:46}
\audiencequestion{Не могли бы вы ещё раз пояснить мысль о возрасте Вселенной и обратной величине Хаббла?}
\lecturerresponse{Для степенных законов получается \(H=p/t\), так что \(H^{-1}\) отслеживает возраст с точностью до множителя порядка единицы. Но в пространстве де Ситтера \(H\) постоянно, тогда как возраст продолжает расти, так что \(H^{-1}\) — это лишь экспоненциальная постоянная времени, а не возраст.}
\end{classroomqa}

\section{Энергия и космический горизонт}

Если плотность энергии вакуума остаётся постоянной, пока Вселенная расширяется, естественно спросить, растёт ли полная энергия вакуума. Ответ, данный здесь, ограничен. В общей теории относительности в учёт энергии входит гравитационная составляющая, так что понятие полной энергии здесь тоньше, чем в элементарной механике. Практически важен момент, связанный с горизонтом.

В точном пространстве де Ситтера космический горизонт имеет фиксированный размер,
\begin{equation}
D=\frac{c}{H}
=\frac{c}{\sqrt{8\pi G\rho_0/3}}.
\end{equation}
Следовательно, объём внутри этого горизонта тоже фиксирован. Энергия вакуума внутри него равна
\begin{equation}
E_{\mathrm{hor}}=\frac{4\pi}{3}D^3\rho_0.
\end{equation}
\editorialnote{Здесь сжато изложен устный расчёт в 00:31:06--00:32:23: берётся размер горизонта, строится объём сферы и выполняется умножение на \(\rho_0\).}
Поскольку и \(D\), и \(\rho_0\) постоянны, эта величина не меняется со временем. Для той части Вселенной, которую данный наблюдатель вообще может когда-либо увидеть, доступная энергия вакуума остаётся постоянной в точном пространстве де Ситтера.

Объекты также могут пересекать космический горизонт примерно в том же смысле, в каком они пересекают горизонт чёрной дыры. Оказавшись снаружи, они выходят из причинной связи с нами.

\begin{classroomqa}
\lecturetimestamp{00:29:27}
\audiencequestion{Если плотность энергии вакуума остаётся постоянной, пока Вселенная расширяется, откуда берётся дополнительная энергия?}
\lecturerresponse{Лекция отвечает на это на языке общей теории относительности, а не в ньютоновских терминах ящика. В точном пространстве де Ситтера размер горизонта остаётся фиксированным, так что энергия вакуума внутри области, которую вообще можно когда-либо увидеть, тоже остаётся фиксированной.}
\end{classroomqa}

\section{Почему была введена инфляция}

Тот же механизм де Ситтера затем переносится в далёкое прошлое, но с гораздо меньшей постоянной времени, порядка
\begin{equation}
H^{-1}\sim 10^{-32}\,\mathrm{s}.
\end{equation}
Такова идея инфляции: ещё до всего, что можно непосредственно наблюдать, Вселенная прошла через период экспоненциального расширения.

Основная мотивация состоит в плоскостности и гладкости Вселенной. Пространственные сечения выглядят очень плоскими. На больших масштабах Вселенная также очень однородна. Ещё поразительнее то, что поверхность последнего рассеяния уже была гладкой, так что это сглаживание не могло произойти только в течение более поздней наблюдаемой истории.

Сдутый воздушный шарик морщинист. Надуйте его достаточно сильно — и любой видимый участок станет гладким. Аналогия с изюмом или черносливом передаёт ту же мысль: морщинистый объект, растянутый в огромное число раз, выглядит гладким на любом видимом участке. Предполагается, что инфляция сделала именно это.

Это также решает проблему горизонта в той форме, в которой она используется здесь. Далёкие друг от друга области поверхности последнего рассеяния обе видны нам, но за более позднюю наблюдаемую историю свет не успел пройти от одной к другой. Поэтому их почти одинаковые плотность и температура указывают на предысторию, в которой одна меньшая область была чрезвычайно растянута.

У инфляции двоякое достоинство. Она устраняет чувствительность к начальным условиям и по той же причине стирает значительную часть информации о том, как началась Вселенная.

Вопрос о плоскостности здесь возникает снова, в более острой форме. Наблюдательно параметр кривизны \(k\) выглядит равным нулю на уровне примерно одного процента или лучше. Геометрия в видимой области выглядит евклидовой. Инфляция объясняет это тем, что делает радиус кривизны намного больше области, которую мы можем видеть.

\begin{classroomqa}
\lecturetimestamp{00:40:05}
\audiencequestion{Предполагаем ли мы равновесие до инфляции?}
\lecturerresponse{Нет. Смысл инфляции почти противоположен этому. Какая бы структура ни существовала до этого, она растягивается настолько сильно, что видимая область теряет память о прежних деталях.}
\end{classroomqa}

\section{Скалярное поле и его потенциал}

Чтобы объяснить, как эффективная энергия вакуума могла быть в ранней Вселенной намного больше, чем сейчас, введём скалярное поле. В отличие от электрического и магнитного полей, которые являются векторными, скалярное поле не имеет направления. Оно просто имеет значение в каждой точке:
\begin{equation}
\phi=\phi(x,y,z,t).
\end{equation}
Это поле гипотетично. Его не отождествляют ни с каким известным в природе полем, а его кванты — ни с какими известными частицами.

Здесь важна та часть его энергии, которая является потенциальной энергией,
\begin{equation}
V(\phi).
\end{equation}
Под ``потенциальной энергией'' понимается энергия, связанная со значением самого поля, а не с производной поля по времени. Разным значениям \(\phi\) могут соответствовать разные энергии.

\begin{figure}[ht]
\centering
\includegraphics[width=0.68\textwidth]{lecture_07_figure_03.png}
\caption{Качественный вид скалярного потенциала \(V(\phi)\), показывающий, что разным значениям поля соответствуют разные потенциальные энергии.}
\label{fig:lecture07-generic-potential}
\lectureframe{00:47:08}
\end{figure}

\begin{classroomqa}
\lecturetimestamp{00:46:34}
\audiencequestion{Что отложено по горизонтальной оси на рисунке \(V(\phi)\)?}
\lecturerresponse{По горизонтальной оси отложено значение поля \(\phi\), а по вертикальной — потенциальная энергия \(V\).}
\end{classroomqa}

Если поле постоянно в пространстве и времени или почти постоянно в рассматриваемой области, то \(V(\phi)\) ведёт себя в точности как энергия вакуума. Оно не разбавляется. Поэтому прежние формулы де Ситтера снова используются после замены \(\rho_0\) на \(V(\phi)\). В частности,
\begin{equation}
H=\sqrt{\frac{8\pi G}{3}V(\phi)}.
\end{equation}

Если \(\phi\) меняется от места к месту, разные области проходят инфляцию с разными характерными временами. Но если растяжение достаточно велико, любой видимый участок становится почти однородным.

\begin{classroomqa}
\lecturetimestamp{01:05:11}
\audiencequestion{Здесь \(\phi\) выступает как параметр времени?}
\lecturerresponse{В этом обсуждении — да, в том ограниченном смысле, что, пока поле скатывается вниз по потенциалу, оно является монотонной функцией времени.}
\end{classroomqa}

\section{Плато, обрыв и повторный разогрев}

Потенциал, необходимый для этой картины, имеет очень плоский, хотя и не идеально плоский, участок, затем доходит до края обрыва, резко падает, не совсем до нуля, и снова возрастает.

\begin{figure}[ht]
\centering
\includegraphics[width=0.68\textwidth]{lecture_07_figure_04.png}
\caption{Инфляционный потенциал с высоким плато, крутым спадом и более низким минимумом.}
\label{fig:lecture07-inflationary-potential}
\lectureframe{01:02:30}
\end{figure}

Эта форма весьма искусственна, но её вводят, потому что космологическая картина, по-видимому, требует чего-то подобного.

Изначально поле находится высоко на плато. Поскольку наклон мал, а эффективная вязкость, связанная с расширением, велика, \(\phi\) меняется лишь очень медленно. Камень, медленно падающий в мёде, ближе к подразумеваемой картине, чем предмет, скользящий по столу. Чем быстрее расширяется Вселенная, тем больше становится эта вязкость.

\begin{classroomqa}
\lecturetimestamp{00:59:07}
\audiencequestion{Что заставляет \(\phi\) двигаться?}
\lecturerresponse{Его толкает наклон потенциала. Наклон мал, а космологическая вязкость велика, поэтому дрейф медленный, а не быстрый.}
\end{classroomqa}

Пока \(\phi\) медленно ползёт по плато, \(V(\phi)\) меняется очень мало. По этой причине Вселенная раздувается почти так, как если бы энергия вакуума была постоянной. Характерное время на этой стадии имеет порядок
\begin{equation}
H^{-1}\sim 10^{-32}\,\mathrm{s}.
\end{equation}
Грубая устная оценка даёт высоту плато около \(10^{-14}\) в планковских единицах.\editorialnote{Это произнесённая оценка ``наилучшей догадки'' на 01:14:38--01:14:47, а не значение, выведенное на лекции.}

Как долго продолжалась инфляция? Это неизвестно, потому что неизвестны ширина и пологость плато. Известно лишь, что размеры Вселенной увеличились в \(e\) по крайней мере примерно 50 раз, и обычно приводят число 60.\editorialnote{Редакторская переформулировка устного обсуждения e-фолдингов на 01:14:05--01:19:27.} Эти величины — e-фолдинги: каждый e-фолдинг — одно умножение на \(e\).

Когда поле достигает края, потенциальная энергия быстро меняется. Поле падает, его потенциальная энергия сначала превращается в кинетическую энергию, а затем — в тепло и частицы, и Вселенная становится горячей. Название ``повторный разогрев'' в этом описании слегка вводит в заблуждение, потому что любая более ранняя тепловая энергия к этому моменту уже была бы разбавлена в ходе инфляции.

Нижний минимум расположен не на нуле. Остающееся там малое значение и связано с нынешним медленным экспоненциальным расширением.

\begin{classroomqa}
\lecturetimestamp{01:19:48}
\audiencequestion{Может ли плато быть в точности плоским, так что возмущения каким-то образом сбрасывали бы поле с края?}
\lecturerresponse{Нет. Оно не может быть совершенно плоским. У него должен быть некоторый наклон. Если бы оно было слишком плоским, то после повторного разогрева Вселенная получилась бы слишком неоднородной.}
\end{classroomqa}

``Большой взрыв'' здесь используется в двух близких смыслах. Он может означать само событие повторного разогрева. Он может также означать самую раннюю непосредственно наблюдаемую стадию, а именно поверхность последнего рассеяния. Неизменный вывод здесь таков: всё, что действительно наблюдается, лежит далеко ниже инфляционного плато.

\begin{classroomqa}
\lecturetimestamp{01:45:42}
\audiencequestion{Является ли сам повторный разогрев Большим взрывом?}
\lecturerresponse{Если угодно, да. Но в этот момент Вселенная всё ещё была чрезвычайно горячей и непрозрачной. Поверхность последнего рассеяния возникает гораздо позже, когда Вселенная достаточно остывает, чтобы стать прозрачной.}
\end{classroomqa}

\section{Квантовые флуктуации и зародыши галактик}

Инфляция делает видимую Вселенную очень однородной, но не идеально однородной. В отсутствие квантовых флуктуаций поле стало бы совершенно гладким. При наличии квантовой механики оно всё же слегка меняется от места к месту.

На центральном рисунке по вертикали отложено обычное пространство, а по горизонтали — значение поля \(\phi\). Если бы поле было точно однородным, то в каждый данный момент каждая точка пространства имела бы одно и то же значение \(\phi\), и все точки вместе дрейфовали бы к краю. Тогда каждая область перевалилась бы через обрыв в один и тот же момент. Квантовые флуктуации нарушают эту точную синхронность.

Суть проясняют две картинки. Одна — это небольшой бугорок в остальном однородном поле. Другая — ряд солдат, идущих под руку к обрыву. Если они выровнены идеально, все падают вместе. Если одни немного впереди, а другие немного позади, то падают они в разное время.

Теперь возникает важный эффект. Область, которая переваливается через край раньше, раньше и разогревается. К тому времени, когда более поздняя область переваливается через край, более ранняя уже успевает несколько расшириться и остыть. Это различие порождает вариации плотности энергии после повторного разогрева.

\begin{classroomqa}
\lecturetimestamp{01:34:00}
\audiencequestion{Более поздно разогретые области горячее или холоднее более ранних?}
\lecturerresponse{Более ранние уже успевают расшириться и остыть. Поэтому более поздние горячее в момент, когда они переваливаются через край, и это различие вносит вклад в неоднородность плотности.}
\end{classroomqa}

Наблюдаемый эффект мал:
\begin{equation}
\frac{\delta\rho}{\rho}\sim 10^{-5},
\end{equation}
примерно одна стотысячная. Это первичные флуктуации плотности, а не тёмная энергия. Это число здесь не выводится из одного лишь качественного наброска; оно используется, чтобы ограничить, насколько высоким и насколько плоским может быть плато.

Они также необходимы. Без них не было бы галактик. После повторного разогрева гравитация усиливает начальный рисунок. Области повышенной плотности притягивают вещество из областей пониженной плотности. Первые становятся ещё более плотными, вторые — ещё менее плотными, и результатом оказывается последующая структура галактик и звёзд.

Это не просто качественная история. Можно вычислить квантовые флуктуации поля в расширяющейся Вселенной, проследить их через повторный разогрев и сравнить результат с картиной флуктуаций, восстановленной по распределению галактик. Согласие очень хорошее.

То же обсуждение приводит и к замечанию об энтропии. Гравитационное сгущение выглядит противоинтуитивно, но второго закона оно не нарушает. В статистической механике гравитация необычна: она стремится вызывать сгущение, и это сгущение увеличивает энтропию, а не уменьшает её.

\begin{classroomqa}
\lecturetimestamp{01:40:24}
\audiencequestion{Как во всём этом проявляется микроволновый фон, учитывая, что \(\phi\) не является электромагнитным полем?}
\lecturerresponse{Первичная картина распределения плотности отпечатывается на поверхности последнего рассеяния. При разогреве энергия поля превращается в обычные частицы и фотоны, а позднее микроволновый фон несёт на себе отпечаток этих первичных флуктуаций.}
\end{classroomqa}

\section{Акустические волны и следующий шаг}

В конце обсуждение ненадолго обращается к акустическим волнам. Они предшествуют последующему гравитационному сгущению. Первичные флуктуации преобразуются в акустические колебания, и эти колебания оставляют на небе неоднородности известного физического размера.

Это важно, потому что известный физический размер на известном расстоянии даёт мерную линейку. Как только измерен угол, под которым он виден на поверхности последнего рассеяния, по геометрии можно определить, является ли пространство плоским, положительно искривлённым или отрицательно искривлённым. Подробности пока откладываются.

\begin{classroomqa}
\lecturetimestamp{01:47:19}
\audiencequestion{А что насчёт акустических волн?}
\lecturerresponse{Они дают мерную линейку на небе. Именно их известный физический размер вместе с их наблюдаемым угловым размером на поверхности последнего рассеяния позволяет измерить геометрию пространства. Подробности переносятся на следующую лекцию.}
\end{classroomqa}
